\completa{11}{Sergio}{Ítem 4: cada dimensión del modelo estrella y cada paso ETL}
\senalar{Comenzar en el hecho central; recorrer las cuatro dimensiones izquierdas y las cuatro derechas. Aclarar que las flechas son relaciones, no una secuencia de carga.}
\decir{El punto cuatro pide llevar los datos a PostgreSQL. De allí nace esta pregunta: ¿cómo almacenamos el dataset conservando el grano y los totales? Grano significa qué representa exactamente una fila. Conservamos el cruce original de tiempo, territorio y perfil; no lo convertimos en una fila por persona. Elegimos un modelo estrella, llamado así porque una tabla central de hechos se relaciona con dimensiones alrededor.}
\decir{El cuadro central dice «hecho desaparicion; PK, id; ocho claves foráneas; mes, cantidad; noventa y tres mil ochocientas treinta y una filas». La tabla de hechos conserva la medida CANTIDAD y las referencias a los catálogos. PK significa clave primaria: id identifica de forma única cada registro almacenado. Ese id es generado en la base; no identifica a una persona del dataset. Una clave foránea referencia una clave válida en otra tabla.}
\decir{A la izquierda, el primer cuadro dice «dim tiempo; PK, anio». Guarda el año y su periodo; su clave es anual, mientras mes queda en el hecho. El segundo dice «dim rango edad; PK, rango edad»: organiza las categorías de edad. El tercero, «dim sexo; PK, sexo», contiene las categorías de sexo. El cuarto, «dim nivel educativo; PK, nivel educativo», describe la categoría educativa. En todos, dim significa dimensión.}
\decir{A la derecha, el primero dice «dim geografia; PK, ubigeo; mil quinientos seis». Relaciona cada código observado con su departamento, provincia, distrito y región. El siguiente dice «dim estado civil; PK, estado civil». Después leemos «dim nacionalidad; PK, nacionalidad». Finalmente, «dim ocupacion; PK, ocupacion». Esos tres catálogos describen estado civil, nacionalidad y ocupación. Ninguna de estas tablas cuenta personas por sí sola: describen las claves utilizadas por el hecho.}
\decir{Cada flecha conecta una dimensión con la tabla central. La relación uno a muchos, escrita uno a N, significa que una fila de catálogo puede ser referenciada por muchos hechos; cada hecho utiliza una clave de ese catálogo. Las flechas no indican que una dimensión se transforme en otra ni el orden de un algoritmo. Hay ocho dimensiones y una tabla de hechos: nueve tablas del modelo, más una tabla intermedia de carga.}
\decir{A la derecha del diagrama aparece el proceso ETL. Paso uno, «Extraer»: el comando barra inversa copy de psql lee el CSV del cliente y carga las noventa y tres mil ochocientas treinta y una filas en stg desaparicion. Stg abrevia staging, una zona intermedia antes del modelo definitivo. Paso dos, «Transformar»: los vacíos categóricos se etiquetan NO ESPECIFICA, los ubigeos se completan a seis dígitos y se obtiene una fila de catálogo por clave. Paso tres, «Cargar»: se insertan primero dimensiones y después hechos con cantidad mayor que cero, dentro de una transacción. La transacción permite confirmar el conjunto de cambios o revertirlo si falla la carga.}
\decir{El control del ubigeo quince cero uno veintiuno es especialmente importante: aparece con dos nombres. Si conserváramos dos filas de catálogo para esa clave, un JOIN podría multiplicar sus cantidades. Agrupar por ubigeo antes de unir evita esa repetición. Finalmente, SUM de cantidad mantiene ciento veintiún mil cuatrocientas veintisiete tanto en la base cargada como en la evidencia de restauración del dump. Un dump es un archivo de respaldo que permite reconstruir la base. Las capturas están enlazadas como evidencia guardada de versión, conteos, esquema, consultas y restauración. Con la base descrita, pasaremos ahora al punto cinco: las consultas Spark SQL.}

